% !TeX root = ../papers/nc-p3.tex
\section{有限片側対称窓の三段加群と電荷四のモナド}
\label{sec:finite-symmetric-instanton-monad}
\begin{theorem}\label{thm:finite-symmetric-instanton-monad}
標準有限Frobenius窓の全辺が有限で、一方の反対交換が対称、反対階数 $\rho>6$ とする。
文字を $U,V,U,Z$ と並べ、$R=R_0\subset U\otimes V$、$K=R^\perp$、
$S=\operatorname{Sym}U^*$ と置く。可換次数加群
\[
 N=(S\otimes V^*)/(S K)
\]
はHilbert列 $(4,6,4)$ の有限長加群であり、$N_2=Z$ という同定を除き、
$w$ は二回の乗法 $S_2\otimes N_0\to N_2$ そのものである。
$N$ の極小自由分解は
\[
 0\to S(-6)^4\to S(-5)^{10}\to
 S(-4)^4\oplus S(-3)^4\oplus S(-2)^4\to S(-1)^{10}\to S^4\to N\to0 .
\]
その最初の線型部分から $\mathbb P(U)$ 上で
\[
 L\otimes\mathcal O(-1)\xrightarrow{a}K\otimes\mathcal O
 \xrightarrow{b}V^*\otimes\mathcal O(1),\qquad
 L=(\textstyle\bigwedge^2 U^*\otimes V^*)\cap(U^*\otimes K),\quad \dim L=4
\]
を得る。$b$ は全射である。$a$ が全点で階数4の開部分ではコホモロジー層は
電荷4の安定な階数2インスタントン束 $\mathcal E$ であり、
$H^1(\mathcal E(n))=N_{n+1}$ は全ての整数で乗法と両立する。
この開部分の窓の添字保持同型は、軸 $U$ を位置0に取ったこれらの束の射影同値と必要十分である。
\end{theorem}
\begin{proof}
$N_0=V^*, N_1=R^*$ であり、$N_2^*$ は三次重なりのうち両端対称な部分である。
$R_1=R^{\rm op}$、全三次重なりは $w(Z^*)$ で既に対称なので $N_2=Z$ を得る。
高次項を調べる。$0\ne t\in N_3^*\subset\operatorname{Sym}^3 U\otimes V$ があれば、
その最小 $U$ 支持を $H$ とする。各 $0\ne\phi\in H^*$ で一回の縮約は非零の $w(z^*)$ である。
この両端の $U$ 支持は有限辺の評価階数から3以上なので $\dim H\ge3$。
$\dim H=3$ なら二回の縮約は全て $P=R\cap(H\otimes V)$ に入り、
$H^\perp$ の点での評価階数が3以上なので $\dim P\le3$。
各 $\phi$ を固定した第二縮約 $H^*\to P$ は階数3で、全非零 $\phi$ で可逆となる。
これは射影平面上の三次行列式の零点で矛盾する。従って $H=U$ で、一回の縮約は
$U^*\simeq N_2^*$ を与える。これで $Z$ を同定すると $w$ は三位置で対称になり、
三つの二対二切断はすべて階数6になってしまう。ゆえに $N_3=0$、以後も零である。

反対の辺 $U,Z$ から作る同じ三段加群は $N$ の次数を反転した体上の双対に等しい。
実際、次数0と2は $Z^*,V$ であり、二乗の作用は同じ $w$ の転置である。
中間の六空間は一段目の作用の核を割って得られ、隣接切断は完全対合なので一致する。
従って両側とも四生成・十本の一次関係のみを持つ。四変数のKoszul分解による
有限長加群とその双対のBetti対称性（外次数 $j$ と $4-j$ の完全対合）により
自由分解の最初と最後の二項が決まる。残る中央項は
$(4+6z+4z^2)(1-z)^4=4-10z+4z^2+4z^3+4z^4-10z^5+4z^6$
から表示の通りとなる。

写像 $b$ は $K$ の十元の一次関係であり、その二次シジジーが $L$ である。
有限長なので層化した $b$ は全射である。$a$ が部分束なら
$\mathcal E=\ker b/\operatorname{im}a$ は階数2で、
$c(\mathcal E)=(1-h)^{-4}(1+h)^{-4}=1+4h^2$ である（$h$ は超平面類）。
十本が独立なので $H^0(\mathcal E)=0$、従って安定である。
実際傾きが非負の階数1部分層があればその反射包 $\mathcal O(d)$ からの非零写像が延び、
$d\ge0$ なので大域切断を与える。二つの短完全列と直線束の中間コホモロジーの消滅により、
$H^1(\mathcal E(n))$ は $K\otimes S_n\to V^*\otimes S_{n+1}$ の余核である。
とくに $H^1(\mathcal E(-2))=0$。これでインスタントンの条件が揃う。
逆にこのコホモロジー加群から $N$、それから二回乗法と $w$ を取り戻せる。
また $N$ からは極小表示の線型部分が基底変更まで一意なので束も戻る。
射影空間の座標変更は $U$ の基底変更、初端と終端の加群の基底変更は $V,Z$ の基底変更である。
軸の同定のスカラー自由度はこれらに含まれる。
有限グラフでの反対同定は射影的には一意である。実際全点を固定する射影変換は、
長さ20の座標四空間を保つ単元倍となり、定理\ref{thm:finite-three-slot-evaluations}の
長さ9以上の安定化単元の証明により定数である。従って余分な印は付かない。
\end{proof}

\begin{corollary}\label{cor:instanton-window-reconstruction}
上の束の集合は次の逆構成で尽くせる。
線型モナド $4\mathcal O(-1)\to10\mathcal O\to4\mathcal O(1)$ のコホモロジー束について
$H^1_*(\mathcal E)$ を取り、$N_0,N_1,N_2$ が指定次元であること、その二回乗法を四窓として
隣接切断・有限Frobenius・有限辺・$\rho>6$ の各有限条件で検査する。
ここではモナドの両端が束として完全、定数列の一次関係は独立とする。
この範囲には全次数正則な窓が実在する。
\end{corollary}
\begin{proof}
モナドから $N$ は既に四生成十一次関係の表示を持っている。$N_1$ が六次元で
二回乗法の隣接階数が6ならその一次関係は $R^\perp$ である。
生成次数2の四シジジーが全ての線型シジジーを張ることは上の自由分解と $a$ の単射から従う。
こうして再構成したモナドも元に戻る。定理\ref{thm:twenty-point-distinct-companion}の
反対階数10の正則窓で、付属計算は $a$ の四次小行列式が四変数の全四次式を張ることを
素数1009で検証する。従って標数零でも共通射影零点はなく、この開部分は空でない。
\end{proof}
